% ------------------------------------------------------------------
% FILE: sections/10_decision_workflow.tex
% Decision workflow for narrowing the algorithm shortlist.
% \input{} into main.tex — do NOT wrap in \documentclass or \begin{document}.
% ------------------------------------------------------------------

\section{Decision Workflow}\label{sec:workflow}

The tiered algorithm strategy (\Cref{sec:algorithms}) replaces the flat 
`funnel'' approach with a \textbf{progressive complexity escalation} workflow. 
Rather than evaluating all candidate models simultaneously to find a single winner, 
the workflow advances through the complexity levels sequentially (Level 0 to Level 3), 
stopping as soon as a tier meets the classification, memory, and latency requirements 
for the target edge hardware.

\paragraph{Design intent.}
This workflow prioritizes the \emph{simplest viable model}. If strong physical separability 
exists, a Level 1 model (Linear SVM) will succeed and offer maximum interpretability 
and minimum deployment cost. Only if non-linear tabular patterns dominate do we escalate 
to Level 2 (Gradient Boosted Trees), and only if complex local spectral correlations 
are necessary do we escalate to Level 3 (1D CNN).

% ------------------------------------------------------------------
\subsection{Workflow Flowchart}
\label{subsec:flowchart}
% ------------------------------------------------------------------

\Cref{fig:decision_workflow} presents the tiered decision workflow as a flowchart.
Process nodes (blue) represent model training and evaluation; decision nodes (orange) 
represent threshold gates; termination/deployment nodes (green) represent success at 
the current tier; and escalation edges (red) represent progression to a higher complexity tier.

\begin{figure}[htbp]
\centering
\begin{tikzpicture}[
  % --- node styles ---
  process/.style={
    rectangle, rounded corners=4pt,
    draw=blue!70!black, fill=blue!8,
    text width=6.5cm, minimum height=1.0cm,
    align=center, font=\small
  },
  decision/.style={
    diamond, aspect=2.0,
    draw=orange!80!black, fill=orange!8,
    text width=3.5cm, minimum height=1.0cm,
    align=center, font=\small
  },
  eliminate/.style={
    rectangle, rounded corners=4pt,
    draw=red!70!black, fill=red!8,
    text width=4.0cm, minimum height=0.8cm,
    align=center, font=\small
  },
  advance/.style={
    rectangle, rounded corners=4pt,
    draw=green!60!black, fill=green!8,
    text width=4.5cm, minimum height=0.8cm,
    align=center, font=\small
  },
  % --- arrow styles ---
  arr/.style={-Stealth, thick},
  yes/.style={arr, green!60!black},
  no/.style={arr, red!70!black},
  plain/.style={arr, blue!70!black},
  node distance=0.8cm and 1.2cm,
]

% --- Nodes ---
\node[process] (l0)
  {\textbf{Level 0: Unsupervised Probe}\\K-Means (=3$) analysis};

\node[decision, below=of l0] (d0)
  {Clusters match\\labels?};

\node[eliminate, right=of d0] (revisit)
  {Revisit Feature\\Representation};

\node[process, below=of d0] (l1)
  {\textbf{Level 1: Linear Baselines}\\Train LR / Linear SVM};

\node[decision, below=of l1] (d1)
  {Meets $, latency,\\\& size limits?};

\node[advance, left=of d1] (deploy1)
  {\textbf{Deploy Level 1}\\Max Interpretability};

\node[process, below=of d1] (l2)
  {\textbf{Level 2: Tabular Ensembles}\\Train Gradient Boosted Trees};

\node[decision, below=of l2] (d2)
  {Meets $, latency,\\\& size limits?};

\node[advance, left=of d2] (deploy2)
  {\textbf{Deploy Level 2}\\Tabular Dominance};

\node[process, below=of d2] (l3)
  {\textbf{Level 3: Deep Spectral}\\Train 1D CNN};

\node[decision, below=of l3] (d3)
  {Meets edge size\\\& latency limits?};

\node[advance, left=of d3] (deploy3)
  {\textbf{Deploy Level 3}\\Spectral Alignment};

\node[eliminate, right=of d3] (fail)
  {Hardware constraints\\exceeded. Abort.};

% --- Edges ---
\draw[plain] (l0) -- (d0);
\draw[no]  (d0.east) -- node[above, font=\footnotesize]{No}  (revisit.west);
\draw[yes] (d0.south) -- node[right, font=\footnotesize]{Yes} (l1.north);

\draw[plain] (l1) -- (d1);
\draw[yes] (d1.west) -- node[above, font=\footnotesize]{Yes} (deploy1.east);
\draw[no]  (d1.south) -- node[right, font=\footnotesize]{No (Escalate)} (l2.north);

\draw[plain] (l2) -- (d2);
\draw[yes] (d2.west) -- node[above, font=\footnotesize]{Yes} (deploy2.east);
\draw[no]  (d2.south) -- node[right, font=\footnotesize]{No (Escalate)} (l3.north);

\draw[plain] (l3) -- (d3);
\draw[yes] (d3.west) -- node[above, font=\footnotesize]{Yes} (deploy3.east);
\draw[no]  (d3.east) -- node[above, font=\footnotesize]{No} (fail.west);

\end{tikzpicture}
\caption{Tiered decision workflow for escalating model complexity. 
         Evaluation halts as soon as a tier meets the quantitative requirements 
         for deployment, avoiding unnecessary complexity.}
\label{fig:decision_workflow}
\end{figure}

% ------------------------------------------------------------------
\subsection{Escalation Gates}
\label{subsec:gate_rationale}
% ------------------------------------------------------------------

\subsubsection*{Level 0 Gate: Fundamental Separability}
The workflow begins with K-Means clustering. If the unlabelled feature space does not 
exhibit structure corresponding to the \texttt{dry}, \texttt{wet}, and \texttt{ice} states, 
no supervised model should be trained. The representation itself must be revisited 
(e.g., investigating different frequency windows or normalization schemes).

\subsubsection*{Level 1 Gate: Linear Sufficiency}
If Linear SVM or Logistic Regression achieve the target macro $ score (e.g., $\geq 0.85$) 
while satisfying edge memory and latency constraints, the workflow terminates. A linear 
model provides absolute transparency regarding which frequency bins drive the decision and 
runs in microseconds on basic microcontrollers. Escalation is only justified if linear 
decision boundaries prove inadequate.

\subsubsection*{Level 2 Gate: Non-Linear Tabular Threshold}
If linear methods fail, Gradient Boosted Trees (GBT) are evaluated. GBTs construct 
non-linear decision boundaries through recursive partitioning and are immune to scaling 
issues. If the GBT ensemble achieves the necessary performance within the hardware budget, 
it is deployed. Escalation to deep learning is reserved for cases where GBTs fail to 
capture correlated spatial/spectral patterns.

\subsubsection*{Level 3 Gate: Spectral Deep Learning}
The 1D CNN is the final tier. It brings the strongest inductive bias (local spectral 
feature correlation) but demands the highest computational cost. If the 1D CNN provides 
the only path to the required $ score, its architecture must be carefully pruned 
and quantised (INT8) to fit within the edge target constraints. If it still exceeds 
the hardware limits, the deployment is aborted.
